% CouncilLens Development Impact Methodology
% Build: tectonic docs/impact-methodology-report.tex
\documentclass[11pt]{article}

\usepackage[margin=1.1in]{geometry}
\usepackage{amsmath, amssymb}
\usepackage{booktabs}
\usepackage{array}
\usepackage{microtype}
\usepackage[hidelinks]{hyperref}
\usepackage{enumitem}
\setlist{itemsep=2pt, topsep=4pt}

\newcommand{\assum}[1]{\texttt{#1}}
\title{CouncilLens Development Impact Methodology}
\author{CouncilLens (CouncilHound) --- impact analysis subsystem}
\date{July 2026 \\ \small model version: joint destination--mode Huff, census-block weights}

\begin{document}
\maketitle

\begin{abstract}
\noindent
CouncilLens produces screening-level impact estimates for proposed local development
projects. A deterministic economic module (spatial-interaction retail capture,
walking activity, an employment ledger) and a deterministic fiscal module
(revenue and service-cost ranges) generate every number; a language model is used
only to extract project facts from public documents under a verbatim-quote
firewall and to write a narrative that a validator constrains to the structured
outputs. This note documents each displayed metric's formula, its academic or
statistical basis, the assumptions that drive its uncertainty range, and the
implementation decisions behind the current system. The estimates rank likely
orders of magnitude under stated assumptions; they are decision-support
evidence, not forecasts.
\end{abstract}

\section{Scope and design principles}

\begin{enumerate}
\item \textbf{Deterministic quantities.} All numbers are produced by closed-form
      models over open data. The language model never supplies a quantity: extracted
      project facts require a verbatim source quote ($\leq$15 words) that is verified
      against the document text in code, and fields without one are nulled. The
      narrative report is post-validated so that every number it contains matches a
      published metric, a stated assumption, a provenance string, or a deterministic
      module note.
\item \textbf{Assumption transparency.} Every contestable parameter is a named
      assumption with a central value, low/high bounds, and a recorded basis. Bounds
      propagate through all formulas by interval arithmetic, so each metric carries a
      range derived from the same assumptions it cites.
\item \textbf{Fail loudly.} Missing inputs (an unpinned tax rate, an absent
      assessment record) degrade the affected metric to ``not computed --- reason,''
      never to a silent default.
\item \textbf{Screening honesty.} Capture, walking, and foot-traffic figures are
      screening estimates in the spatial-interaction tradition; the report states
      this in plain language and names the most sensitivity-driving assumptions.
\end{enumerate}

\section{Notation and data inputs}

\begin{tabular}{ll}
$U$        & proposed dwelling units (extracted from project documents) \\
$\rho$     & occupancy rate at stabilization (\assum{occupancy\_rate}) \\
$\bar{s}$  & average renter household size (ACS B25010, population-weighted) \\
$y$        & median household income of the site's census tract (ACS B19013) \\
$\pi$      & new-construction income premium (\assum{income\_premium\_new\_construction}) \\
$e_c$      & CES average annual expenditure for category $c$ (BLS CE 2023) \\
$y_{\mathrm{CES}}$ & CES average pre-tax income per consumer unit (\$101{,}805, 2023) \\
$\kappa$   & CES scaling factor (\assum{ces\_scale}) \\
$A_j$      & attractiveness of destination $j$ (POI count~$=1$ per business) \\
$t^w_j,\ t^d_j$ & walk / drive travel time (minutes) from the site to $j$ \\
$\beta_w,\ \beta_d$ & walk / drive impedance decay per minute \\
$w_c$      & zero-impedance walk preference for category $c$ \\
$\tau$     & daily walk trips per resident (\assum{walk\_trips\_per\_resident\_day}) \\
\end{tabular}

\medskip
\noindent
Networks are OpenStreetMap walk and drive graphs built with OSMnx
\cite{boeing2017}; walk edge times use 4.8\,km/h. Population weights use 2020
Census (PL~94-171) block centroids. Points of interest merge Overture Maps,
OpenStreetMap, and (when licensed) Foursquare OS Places into a nine-class
taxonomy. Travel times come from single-source Dijkstra runs over edge minutes;
each business inherits the time of its nearest network node.

\section{Demand module}

\subsection{New households and residents}
\begin{align}
N_{hh}  &= U \cdot \rho \\
N_{res} &= N_{hh} \cdot \bar{s}
\end{align}
Occupancy at stabilization is benchmarked against the Census Housing Vacancies
and Homeownership series \cite{hvs}; $\bar{s}$ is estimated from ACS B25010
(renter-occupied) over the city's populated block groups with population
weighting, carried with a $\pm 0.3$ sensitivity band. \emph{Limitation:} this
tenure-level average is blind to unit mix; bedroom-mix demographic multipliers
in the tradition of Listokin et al.~\cite{listokin2006} would lower resident
estimates for studio/one-bedroom-heavy multifamily and are the recommended
refinement (\S\ref{sec:limitations}).

\subsection{Aggregate income and category spending}
\begin{align}
Y   &= N_{hh} \cdot y \cdot \pi \\
S_c &= Y \cdot \frac{e_c}{y_{\mathrm{CES}}} \cdot \kappa
\end{align}
Tract selection follows the site polygon; when tract medians are suppressed the
model falls back to the citywide population-weighted mean and records the
substitution. Category shares come from the BLS Consumer Expenditure Survey
2023 tables \cite{blsces}, income-scaled linearly from the CES average
consumer unit. Linear scaling overstates necessity spending at high incomes
(Engel effects visible in the CE income-quintile tables); $\kappa$'s
$\pm 15\%$ band partially absorbs this (\S\ref{sec:limitations}).

\section{Retail capture: joint destination--mode choice}
\label{sec:capture}

Retail capture follows the probabilistic trade-area tradition of Huff
\cite{huff1963, huff1964} with the exponential impedance of the
entropy-maximizing spatial-interaction family \cite{wilson1971}. Destination
and mode are chosen \emph{jointly}, as a multinomial logit \cite{mcfadden1974,
benakiva1985} over (business, mode) pairs. For category $c$, business $j$, and
mode $m \in \{w, d\}$:
\begin{equation}
P_c(j, m) \;=\;
\frac{A_j^{\alpha}\, \omega_{c,m}\, e^{-\beta_m t^m_j}}
     {\sum_k A_k^{\alpha} \left( \omega_{c,w}\, e^{-\beta_w t^w_k}
       + \omega_{c,d}\, e^{-\beta_d t^d_k} \right)},
\qquad
\omega_{c,w} = w_c,\ \ \omega_{c,d} = 1 - w_c,
\label{eq:joint}
\end{equation}
with $\alpha = 1$. The mode preference $w_c$ has a clean interpretation: at
zero travel time for both modes, the walk share is exactly $w_c$; as walk time
grows, walking loses utility to driving and the trip (and its dollars) arrives
by car instead. Annual capture and walk-arriving capture at business $j$:
\begin{equation}
C_j = \sum_c S_c \sum_m P_c(j, m), \qquad
C^{walk}_j = \sum_c S_c \, P_c(j, w).
\end{equation}
Two invariants are pinned by unit tests: $\sum_j C_j = \sum_c S_c$ (joint
probabilities sum to one), and $\sum_j C^{walk}_j \leq \sum_c S_c w_c$ with
equality only when every destination is at zero travel time --- the walk total
is \emph{endogenous}, not an assumption restated.

\paragraph{Own-retail destination.} The project's ground floor competes inside
the model rather than being appended outside it. Its retail square footage is
converted to POI-count-equivalent attractiveness
($A_{own} = \mathrm{sqft} / \assum{own\_retail\_sqft\_per\_equiv\_poi}$),
distributed across a fixed category mix (50\% restaurant/bar, 15\% comparison,
15\% convenience, 20\% personal services), at zero travel time from the site.
Huff's original attractiveness variable was selling area \cite{huff1963};
competitor attractiveness here is POI count because open POI data lacks floor
area, so the own-retail conversion reconciles the two units.

\paragraph{In-city shares.} For food-away and for all-retail, the in-city
capture share is the probability mass at destinations inside the city boundary,
$\theta = \sum_{j \in \mathrm{city}} \sum_m P(j,m)$. These shares are the
cross-module inputs to the meals-tax and sales-tax lines.

\paragraph{Reporting rollup.} Businesses are grouped for display by DBSCAN
\cite{ester1996} ($\varepsilon = 150$\,m, $\mathrm{minPts} = 3$) and each
cluster is labeled with the name of the member business nearest its centroid;
cluster values are sums of member-business captures and play no role in the
allocation itself.

\section{Walking activity}

\subsection{Resident walk trips and exact marginal flows}
Daily walk trips are $T = N_{res} \cdot \tau$. These all-purpose trips are
allocated over POI-weighted walk-network nodes with exponential decay and
routed along shortest paths --- deterministically, with no sampling --- giving
per-street-segment flows:
\begin{equation}
f_e \;=\; \sum_{v} T \cdot
\frac{g_v\, e^{-\beta_w t_v}}{\sum_u g_u\, e^{-\beta_w t_u}}
\cdot \mathbb{1}\!\left[ e \in \mathrm{SP}(\mathrm{site} \to v) \right],
\end{equation}
where $g_v$ is the POI weight at node $v$ and $\mathrm{SP}$ is the shortest
path. The change caused by the project \emph{is} the trips its residents
generate, so no baseline differencing (and no sampling noise) enters the
numerator.

\subsection{Baseline index and percent comparison}
The comparison baseline is a seeded, sampled, weighted betweenness index
\cite{freeman1977, brandes2001} of today's walkers: $k = 2000$
origin--destination pairs drawn with origins proportional to census-block
population within 400\,m of each node and destinations proportional to POI
weight, routed by shortest paths. Scaling the sampled counts by
$(\mathrm{total\ population\ weight} \times \tau) / k$ expresses the baseline
in the same trips units, and the trip rate $\tau$ cancels from the percentage:
\begin{equation}
\Delta_e \;=\; 100 \cdot
\frac{f_e}{n_e \cdot \bar{P}\, \tau / k}
\quad \text{(reported only where } n_e \geq 5\text{)},
\end{equation}
with $n_e$ the sampled baseline count on edge $e$ and $\bar{P}$ the total
population weight. The headline metric averages $\Delta_e$ over the ten
commercial street segments nearest the site. Without calibrated pedestrian
counts this is an index comparison, not a volume estimate; the correlation of
betweenness-type indices with observed pedestrian volumes is supported by the
network-analysis literature but is context-dependent.

\subsection{Consistency diagnostic}
\begin{equation}
\text{spending per walk trip} = \frac{\sum_j C^{walk}_j}{365\, T}.
\end{equation}
Because $T$ counts \emph{all} walking (exercise, transit access, dog walks)
while the numerator is shopping only, a plausible value is a few dollars per
average trip (implying \$5--\$15 per actual shopping trip). Values outside
\$0.25--\$30 add a consistency flag to the report.

\section{Employment ledger}

\begin{equation}
J^- = \frac{\mathrm{existing\ commercial\ sqft}}{\assum{sqft\_per\_office\_job}},
\qquad
J^+ = \frac{\mathrm{proposed\ retail\ sqft}}{\assum{sqft\_per\_retail\_job}},
\qquad
\Delta J = J^+ - J^-.
\end{equation}
Space-per-worker values (300 [200--450] office; 500 [400--700] retail, gross)
sit between dense workplace-utilization benchmarks and the building-stock
averages observable in CBECS \cite{cbecs}. These are site-activity screens
only; no indirect or induced employment is modeled.

\section{Fiscal module}

The module reports both of the two canonical fiscal-impact framings ---
average(per-capita) cost and marginal cost --- following Burchell \&
Listokin's classification \cite{burchell1978} and the planning literature's
standard caution that neither is universally correct for infill
\cite{kotval2006}.

\subsection{Revenue}
\begin{align}
\text{current RE tax} &= AV_{now} \cdot r / 100 \\
AV_{proj} &= \mathrm{median}\!\left(\tfrac{AV_i}{\mathrm{units}_i}\right)
             \cdot U \;+\; \mathrm{retail\ sqft} \cdot p_{sqft}
             \qquad \text{(bounds: comp 25th--75th pct)} \\
\text{projected RE tax} &= AV_{proj} \cdot r / 100 \\
\text{RE tax increase} &= (AV_{proj} - AV_{now}) \cdot r / 100 \\
\text{meals tax} &= S_{\mathrm{food\ away}} \cdot \theta_{\mathrm{food}} \cdot r_{\mathrm{meals}} \\
\text{local sales tax} &= \left( S_{\mathrm{grocery}} + S_{\mathrm{comparison}}
                          + S_{\mathrm{convenience}} \right) \cdot \theta_{\mathrm{all}}
                          \cdot r_{\mathrm{sales}} \\
\text{personal property} &= N_{hh} \cdot p_{pp}
    \qquad \text{(not computed until $p_{pp}$ is pinned)}
\end{align}
Comparables are apartment-class parcels (land-use code 352) built within a
15-year lookback, pulled from the city's public assessment database with
per-unit values $AV_i / \mathrm{units}_i$; fewer than three comps degrades the
metric with a low-confidence flag. All rates ($r$, $r_{\mathrm{meals}}$,
$r_{\mathrm{sales}}$) are jurisdiction configuration pinned with source URL and
fiscal year. Virginia's 1\% local-option sales tax applies to food for home
consumption, so grocery is correctly in the taxable base; personal services are
excluded as generally non-taxable. Business-license (BPOL) revenue is not
computed in this version.

\subsection{Service costs and net impact}
\begin{align}
\text{naive per-capita cost} &= N_{res} \cdot \frac{GF}{P} \\
\text{marginal cost} &= \text{naive} \cdot \mu, \qquad
\mu = \assum{marginal\_cost\_factor} \\
\text{net}_{\,method} &= R_{inc} - \text{cost}_{\,method}
\end{align}
where $GF/P$ is General Fund expenditure per capita and $R_{inc}$ is
\emph{incremental} recurring revenue (the real-estate component is the tax
\emph{increase} when the current tax is known). The naive method is an
intentional upper bound: it allocates fixed citywide costs to incremental
residents. The marginal framing scales only the cost share expected to grow
and can understate costs if capacity thresholds are crossed. The headline range
spans both framings; per-method nets are also published. School impact is
reported as a note --- $U \times \assum{students\_per\_unit}$ students, per
multifamily generation-rate evidence \cite{listokin2006} --- and deliberately
does not enter the arithmetic of either cost method (per-capita costing already
embeds average school costs).

\section{Assumptions}

Central values with propagated bounds; every metric lists the assumption keys
that touched it.

\begin{center}
\small
\begin{tabular}{@{}llll@{}}
\toprule
Key & Value & Range & Basis \\
\midrule
\assum{occupancy\_rate} & 0.95 & 0.90--0.97 & stabilized MF occupancy; cf.\ HVS vacancy \cite{hvs} \\
\assum{avg\_hh\_size\_renter} & ACS-derived & $\pm 0.3$ & ACS B25010 renter-occupied, pop-weighted \\
\assum{income\_premium\_new\_construction} & 1.125 & 1.0--1.25 & screening premium, stated openly \\
\assum{ces\_scale} & 1.0 & 0.85--1.15 & CES national-average drift allowance \cite{blsces} \\
\assum{beta\_walk} (min$^{-1}$) & 0.15 & 0.07--0.15 & impedance decay; set to top of band (2026-07) \\
$\beta_d$ (min$^{-1}$, fixed) & 0.15 & --- & drive impedance decay \\
\assum{walk\_share\_neighborhood} & 0.60 & 0.40--0.80 & zero-impedance walk preference \\
\assum{walk\_share\_comparison} & 0.10 & 0.00--0.30 & comparison goods travel by car \\
\assum{walk\_share\_grocery\_entertainment} & 0.30 & 0.10--0.50 & interpolated \\
\assum{walk\_trips\_per\_resident\_day} & 2.0 & 1.0--3.0 & see calibration flag, \S\ref{sec:limitations} \\
\assum{sqft\_per\_office\_job} & 300 & 200--450 & planning standard; cf.\ CBECS \cite{cbecs} \\
\assum{sqft\_per\_retail\_job} & 500 & 400--700 & planning standard; cf.\ CBECS \cite{cbecs} \\
\assum{own\_retail\_sqft\_per\_equiv\_poi} & 2{,}000 & 1{,}500--3{,}000 & typical inline suite size \\
\assum{commercial\_value\_per\_sqft} & \$275 & \$200--\$400 & screening; replace with comps \\
\assum{students\_per\_unit} & 0.175 & 0.10--0.25 & MF generation rates \cite{listokin2006} \\
\assum{marginal\_cost\_factor} & 0.40 & 0.25--0.60 & share of average cost that scales \\
\bottomrule
\end{tabular}
\end{center}

\noindent
Structural constants: walk speed 4.8\,km/h; DBSCAN $\varepsilon=150$\,m,
$\mathrm{minPts}=3$; network extent = city boundary + 3\,km; POI study area =
boundary + 2\,mi; node population radius 400\,m; baseline sampling $k = 2000$,
fixed seed; comp lookback 15 years.

\section{Implementation decisions}

\begin{itemize}
\item \textbf{Per-business destinations.} The Huff run treats every retail POI
      as an individual destination (a single-source Dijkstra makes per-POI times
      free); DBSCAN clusters exist only as reporting rollups. Capture is therefore
      spatially resolved to business locations, and map heat layers consume real
      per-business weights rather than interpolated cluster centroids.
\item \textbf{Joint rather than nested choice.} Equation~\eqref{eq:joint} is a
      single multinomial logit over (destination, mode) pairs. A nested logit would
      relax the independence-of-irrelevant-alternatives property across modes; the
      joint form was chosen for transparency and because screening bounds on $w_c$
      and $\beta_w$ dominate the error budget.
\item \textbf{Exact marginal flows.} Project-caused foot traffic is computed as
      the shortest-path routing of the residents' own trips, replacing an earlier
      sampled before/after betweenness difference whose noise exceeded the signal at
      this project scale. Sampling now appears only in the baseline denominator.
\item \textbf{Census-block population.} Node population weights use 2020 census
      block centroids (the finest published geography) fetched by spatial envelope,
      which also extends the baseline beyond the city line that the county-scoped
      ACS layer stopped at.
\item \textbf{Determinism.} All sampling is seeded; two evaluations on a warm
      cache produce identical metric values.
\item \textbf{LLM firewall.} Document extraction demands a verbatim
      $\leq$15-word quote per numeric field, verified in code against the document
      text; unverifiable values are nulled. Parcel identifiers are verified the
      same way.
\item \textbf{Narrative validator.} The generated report is scanned for
      numbers; each must match a metric value/bound, spec quantity, assumption, or
      deterministic note within rounding, else the draft is regenerated once and
      then hard-fails.
\item \textbf{Assessment data.} The city publishes no bulk assessment layer, so
      site values and multifamily comparables come from the public Patriot WebPro
      search interface, throttled to $\sim$1 request/s and cached; owner names are
      neither parsed nor stored.
\item \textbf{Map scales.} Heat layers are clipped to commercial-retail zoning
      polygons and colored on dollar-pinned viridis ramps (nice-number ceilings),
      so a color means the same dollars regardless of layer toggling; capture and
      walk-in capture share one scale to make their comparison meaningful.
\end{itemize}

\section{Known limitations and audit notes}
\label{sec:limitations}

\begin{enumerate}
\item \textbf{Household size is unit-mix-blind.} ACS renter-average
      $\bar{s} \approx 3.1$ for this city likely overstates residents of a
      studio/1BR/2BR multifamily building; bedroom-mix multipliers
      \cite{listokin2006} typically imply 1.5--2.1. Resident-scaled quantities
      (spending, walk trips, per-capita costs) inherit this bias; the naive-cost
      side inflates faster than revenue, making the fiscal range conservative.
\item \textbf{Median $\times$ households $\neq$ aggregate income.} Using tract
      B19013 understates aggregate income under right-skewed distributions; ACS
      B19025 (aggregate income) or mean income would be the better estimator.
\item \textbf{Linear CES scaling.} Category spending scales linearly with
      income; Engel effects mean necessity categories are overstated at high
      incomes. Income-quintile CE tables would refine this.
\item \textbf{Walk-trip rate calibration flag.} $\tau = 2.0$ trips/resident/day
      exceeds national NHTS person-level walking rates (roughly 0.3--0.7 in
      metropolitan samples \cite{nhts}); walkable mixed-use contexts run higher,
      but 2.0 should be treated as an upper-range choice pending local evidence.
      The foot-traffic \emph{percentage} is unaffected ($\tau$ cancels), and
      walk-in dollars are independent of $\tau$; the trips metric and street-flow
      magnitudes scale with it.
\item \textbf{IIA.} The joint logit imposes proportional substitution across all
      (destination, mode) pairs; a nested structure would be more standard.
\item \textbf{One-sided $\beta_w$ band.} The central value sits at the top of
      its sensitivity range (analyst calibration), so sensitivity sweeps only
      explore gentler decay.
\item \textbf{Attractiveness proxy.} Competitor attractiveness is POI count
      (floor area being unavailable in open POI data), while own-retail
      attractiveness is derived from square footage; the conversion constant
      carries the reconciliation and its bounds.
\item \textbf{Sales-tax base share.} The taxable base applies the all-retail
      in-city share to the three taxable categories; a category-specific share
      would be marginally more precise.
\item \textbf{Foot-traffic index is uncalibrated.} Percent changes compare
      modeled flows to a sampled opportunity index, not to counted pedestrians;
      local counts would convert the layer to volumes.
\item \textbf{Fiscal method dependence.} The net range spans average-cost and
      marginal-cost framings on purpose; the midpoint of that range is a
      presentation convenience, not an estimate with independent standing.
\end{enumerate}

\section{Data sources}

City of Fairfax GeoHub ArcGIS services (boundary, parcels, zoning); City of
Fairfax public assessment database (Patriot WebPro); City budget pages (tax
rates, General Fund, pinned with fiscal year); Census ACS 5-year (B01003,
B19013, B25010, B25044, B08301) and 2020 PL~94-171 block population via
TIGERweb; LEHD LODES~8; OpenStreetMap via OSMnx \cite{boeing2017}; Overture
Maps places; BLS Consumer Expenditure Survey 2023 \cite{blsces}; FHWA National
Household Travel Survey \cite{nhts}; Virginia Code \S 58.1-605 (local-option
sales tax).

\begin{thebibliography}{99}\small

\bibitem{huff1963} Huff, D.~L. (1963). A probabilistic analysis of shopping
center trade areas. \emph{Land Economics}, 39(1), 81--90.

\bibitem{huff1964} Huff, D.~L. (1964). Defining and estimating a trading area.
\emph{Journal of Marketing}, 28(3), 34--38.

\bibitem{wilson1971} Wilson, A.~G. (1971). A family of spatial interaction
models, and associated developments. \emph{Environment and Planning A}, 3(1),
1--32.

\bibitem{mcfadden1974} McFadden, D. (1974). Conditional logit analysis of
qualitative choice behavior. In P.~Zarembka (Ed.), \emph{Frontiers in
Econometrics} (pp.~105--142). Academic Press.

\bibitem{benakiva1985} Ben-Akiva, M., \& Lerman, S.~R. (1985). \emph{Discrete
Choice Analysis: Theory and Application to Travel Demand}. MIT Press.

\bibitem{freeman1977} Freeman, L.~C. (1977). A set of measures of centrality
based on betweenness. \emph{Sociometry}, 40(1), 35--41.

\bibitem{brandes2001} Brandes, U. (2001). A faster algorithm for betweenness
centrality. \emph{Journal of Mathematical Sociology}, 25(2), 163--177.

\bibitem{ester1996} Ester, M., Kriegel, H.-P., Sander, J., \& Xu, X. (1996). A
density-based algorithm for discovering clusters in large spatial databases
with noise. \emph{Proceedings of KDD-96}, 226--231.

\bibitem{boeing2017} Boeing, G. (2017). OSMnx: A Python package to work with
graph-theoretic OpenStreetMap street networks. \emph{Journal of Open Source
Software}, 2(12), 215.

\bibitem{burchell1978} Burchell, R.~W., \& Listokin, D. (1978). \emph{The
Fiscal Impact Handbook: Estimating Local Costs and Revenues of Land
Development}. Center for Urban Policy Research, Rutgers University.

\bibitem{kotval2006} Kotval, Z., \& Mullin, J. (2006). \emph{Fiscal Impact
Analysis: Methods, Cases, and Intellectual Debate}. Lincoln Institute of Land
Policy working paper WP06ZK2.

\bibitem{listokin2006} Listokin, D., Voicu, I., Dolphin, W., \& Camp, M.
(2006). \emph{Who Lives in New Jersey Housing? Updated New Jersey Demographic
Multipliers}. Center for Urban Policy Research, Rutgers University.

\bibitem{blsces} U.S. Bureau of Labor Statistics. \emph{Consumer Expenditure
Surveys}, 2023 tables. \url{https://www.bls.gov/cex/}

\bibitem{nhts} Federal Highway Administration. \emph{National Household Travel
Survey}, 2022. \url{https://nhts.ornl.gov/}

\bibitem{hvs} U.S. Census Bureau. \emph{Housing Vacancies and Homeownership
(CPS/HVS)}. \url{https://www.census.gov/housing/hvs/}

\bibitem{cbecs} U.S. Energy Information Administration. \emph{Commercial
Buildings Energy Consumption Survey (CBECS)}, 2018 building characteristics
tables. \url{https://www.eia.gov/consumption/commercial/}

\bibitem{ewing2010} Ewing, R., \& Cervero, R. (2010). Travel and the built
environment: A meta-analysis. \emph{Journal of the American Planning
Association}, 76(3), 265--294.

\end{thebibliography}

\end{document}
